\section{Discussion}
\label{sec:discussion}

\subsection{Relation to prior monitoring work}

AWD monitoring has been approached from three directions. \emph{Instrumented}: a
low-cost wireless water-level sensor for lowland rice under
AWD~\cite{delacruz2022}, which solves the measurement problem but multiplies
device cost per plot. \emph{Remote sensing}: identification of AWD adoption in
the Vietnamese Mekong Delta~\cite{lovell2019}, water-status monitoring with
ALOS-2 scattering-power decomposition~\cite{islam2026alose}, and region-scale
machine-learning mapping of water-saving practices~\cite{mlrice2025}; these cover
large areas cheaply but revisit every two to four days and cannot resolve the
$72$\,h dry phase that determines credit eligibility. \emph{Programme level}: the
one-million-hectare high-quality low-emission rice programme~\cite{qd1490}, which
had brought over $350\,000$ hectares under cultivation by 2026~\cite{mae2026progress},
and MRV identified as the missing link for low-emission rice in Southeast
Asia~\cite{irri2026mrv}.

On the agronomic side, meta-analysis quantifies the water and yield effects of
AWD~\cite{carrijo2017}, while adoption economics~\cite{lampayan2015} and the
technical and institutional barriers to scaling~\cite{enriquez2021} explain why
instrumented monitoring has not spread. None of these lines evaluates the
\emph{reliability of the action record itself}, which is the question asked here.

The gateway supplying the device-anchored timestamp is the controller of the
companion work~\cite{twingate2026}, whose command dispatch runs over a
burst-noise channel of the Gilbert--Elliott type~\cite{gilbert1960,elliott1963}.
Our own prior scheduling results for safety-critical agricultural IoT networks
address the same transmission regime~\cite{rabs,cawvou}, but neither establishes
an evidence record, which is the function relied on here. Because that controller
already signs every command it issues, the anchor adds storage and transmission
of a record that exists regardless --- which is what keeps the marginal cost of
attestation near zero.

\subsection{What survives of the original claim}

The premise that motivated this work was attractive: an irrigation command is
already a digital event, so MRV evidence should be nearly free. That premise
survives, but not in the form we first published. It survives because the
controller that issues the command can also sign it, and the signed record is
what makes the evidence admissible --- not because a hydrological model can
detect a lie. Where the original draft claimed that physics attests the log, the
correct statement is that physics \emph{constrains} the log: it rules out depths
no soil can hold and states no field can reach, which is worth having, and it is
silent about the one transformation that pays.

\subsection{Why publishing the threshold is not the problem}

One might conclude that $\varepsilon$ should be kept secret. It should not, for
two reasons. A verifier must be able to reproduce a decision, which requires
published parameters. And secrecy would not help: by Proposition~\ref{prop:blind}
the time-shifting attack preserves the total volume and every individual depth,
so it does not need to know $\varepsilon$ in order to stay beneath the
infiltration bound --- it needs only to avoid moving water, which is visible to
any farmer. The vulnerability is structural, not a consequence of disclosure.

\subsection{Relation to remote sensing}

Satellite radar is the obvious low-cost alternative, so rather than argue about it
we measured it. Table~\ref{tab:benchmark} scores every route by one quantity: how
many of the $150$ credit-eligible dry phases actually present across the $60$
simulated fields it could confirm. A phase counts as confirmable only if at least
one observation falls inside it, because a $72$\,h phase lying wholly between two
passes cannot be evidenced at all --- which costs the farmer a credit even when
the practice was performed correctly. Single-satellite Sentinel-1, revisiting every
twelve days~\cite{clauss2017s1}, confirms $90/150$ ($60\%$); the full A+B
constellation at six days~\cite{hashemi2022s1} reaches $124/150$ ($83\%$), and
still misses one phase in six at zero device cost. A continuous water-level tube
confirms all $150$, at one device per plot~\cite{delacruz2022}. The signed command
record also confirms all $150$ --- because the phases are defined by the intervals
between actuations, which the controller already timestamps --- at no additional
device, since the controller is the one that opened the pump.
Two cautions belong with that comparison. The confirmation column measures the
ability to \emph{earn} a credit, not to detect fraud: the tube and the command
record tie at $100\%$ here and differ entirely in cost and in what they can prove
about a doctored log. And photograph-plus-self-report scores $0\%$ not because it
is worthless but because it produces no observation series to reconcile against;
every confirmation under it must come from a field visit.
Radar therefore keeps the role VM0051 Appendix 4 assigns it: a per-plot cross-check
that lowers the frequency of field verification~\cite{vm0051}.

\subsection{Limitations}

\begin{enumerate}
\item \emph{Simulated farmers.} Forgeries are injected into simulated reference
trajectories. A field pilot is needed to measure behavioural response, in
particular whether farmers shift their \emph{actual} pumping once the anchor is
known to exist --- an attack our model cannot represent, since it changes $M$
rather than $L$.
\item \emph{The anchor assumes an honest gateway.}
Proposition~\ref{prop:complete} is completeness against log forgery, not against
a compromised controller. Device integrity, key custody and replay protection
are assumed from the companion work and are not re-established here.
\item \emph{Not applicable to manual irrigation.} Where a sluice is opened by
hand there is no command record, so the anchor is unavailable. This excludes a
substantial share of delta farms and bounds the addressable population of the
scheme.
\item \emph{The attack is demonstrated in the water-scarce year only.} In the
water-surplus year (2025) farmers need so little irrigation that most simulated
logs fall below the three-command minimum at which a timing attack is even
constructible: of $37$ genuinely altered logs under the full $\Theta$, all $37$
come from 2024. The escarpment between the two years is itself a finding --- MRV
fraud risk is concentrated where water is scarce --- but it also means the
attack results are not yet evidence about a wet season. Under the reduced
$\Theta$ twelve altered logs do come from 2025, and they behave the same way,
which is suggestive rather than conclusive.
\item \emph{$\Theta$ is a judgement.} Five textures, two $K_c$ regimes, three
efficiencies and five initial states bracket the delta but do not exhaust it. A
field outside $\Theta$ loses the guarantee of Proposition~\ref{prop:noalarm},
which is exactly the situation the proportional threshold and the query verdict
exist to absorb.
\item \emph{Device prices are planning estimates, not quotations.} The tier
ordering and the break-even inequality are unaffected; the absolute margin is.
\item \emph{No emission quantification.} The layer produces activity data.
Coupling it to a methane model and locally calibrated emission factors is
separate work, as is the cost table's assumption that credits scale with
$n_{\mathrm{dry}}$.
\end{enumerate}

\section{Conclusion}

We set out to make rice MRV nearly free by auditing a farmer-written irrigation
log against soil-water physics, and the first version of that layer reported no
false accusation and complete detection. Both numbers were artefacts of a field
generated by the auditing model itself. Rebuilding the field along seven axes of
realism, and replacing fixed forgery recipes with an adversary that optimises the
credited quantity, changes the conclusion.

The point-model detector falsely accuses $26.7\%$ of honest farmers, because an
absolute pump-gap threshold cannot distinguish a forgotten entry from a secret
irrigation. A proportional threshold and an intersection over a $150$-element
uncertainty set reduce that to $5.0\%$ and $0/60$ respectively, at a measured
cost in detection of omission fraud, from $9/9$ to $4/9$ on the logs that were
actually altered. Meanwhile a volume-preserving, depth-bounded permutation of
claimed times inflates creditable dry phases by $33\%$ to $38\%$ in paired
comparison while escaping the robust physical tier in $8$ of $9$ and $3$ of $3$
altered logs --- and the point model that catches more of it is the same model
that accuses $26.7\%$ of honest farmers, so no physical configuration escapes
the trade-off. We prove this is unavoidable for any detector whose verdict
is a function of the replayed trajectory: delaying a claim both lengthens the dry
phase and increases the deficit available to it, so profit and undetectability
share a gradient.

The repair is not hydrological. Matching the claimed log against the gateway's
signed, time-stamped command record catches $116/116$ genuinely forged logs and
accuses none of $60$ honest ones, and it costs nothing beyond a controller the
companion work already deploys. The honest formulation of the result is
therefore narrower and more useful than the one we started with: near-zero-cost
MRV for low-emission rice is achievable, but only where a device stands between
the farmer and the pump. Three previously published numbers are corrected
accordingly --- the admissible depth is the range $[49.0,86.0]$\,mm rather than a
single $61$\,mm; the third physical filter cannot activate on dry-season delta
rice, clipped depletion peaking at $0.9987$ over $150$ configurations; and a
subthreshold-based absolute gap test is inadmissible. Code, both uncertainty
sets, all $281$ audited logs and $46$ unit tests, one per design hole, are
public.
